\subsection{疲劳分析}
\label{sec:fatigue-mag}

\subsubsection{取值}

大容量强磁吸附工具承力件材料同为 4330V-155KSI，所用 S-N 曲线与第~\ref{sec:fatigue-4in1}~节相同：$1\times10^{6}$ 次处 $S_e=191.5$\,MPa，$1\times10^{3}$ 次处许用应力幅 929.7\,MPa，$\sigma_a=4\,512.7\,N^{-0.2287}$。有限元模型、接触设置与网格沿用本章强度校核，额定工况为轴向拉力 1\,000\,kN 与扭矩 17\,kN$\cdot$m。

轴向拉伸工况的峰值 874.7\,MPa 位于上接头与心轴连接接头的螺纹根部，系几何尖角与键合面终止线交汇处的数值奇异；距该点约 3\,mm、梯度已趋平缓处的评定应力为 604.1\,MPa（节点 673941），相邻节点为 594.9\,MPa，二者相差 1.5\%。

抗扭工况按完全反向 $R=-1$ 取定。此处与轴向拉伸取法不同，原因在于载荷而非结构：扭矩在正反转、倒划眼与憋钻释放过程中会换向，完全反向可覆盖最不利情形，此时平均应力为零、应力幅等于应力峰值，较脉动循环严格一倍。峰值 979.7\,MPa 位于同一接头的螺纹根部，评定应力取 597.65\,MPa（节点 673622）。同一位置更靠近奇异点的节点读数为 774.29\,MPa，两者相差 30\%，表明该处应力梯度陡峭；按统一规则取梯度平缓区的 597.65\,MPa，774.29\,MPa 单列备查。评定结果列于表~\ref{tab:fatigue_mag}。

\begin{table}[htbp]
\centering
\caption{大容量强磁吸附工具的疲劳评定结果（200 趟）}
\label{tab:fatigue_mag}
\begin{tabularx}{\textwidth}{lCCC}
\hline
项目 & 轴向拉伸（$R=0$） & 抗扭（$R=-1$） & 判据 \\
\hline
额定工况 & 1\,000\,kN & 17\,kN$\cdot$m & --- \\
评定应力 / MPa & 604.1 & 597.7 & --- \\
应力幅与平均应力 / MPa & 302.1 & 597.7 / 0 & --- \\
等效对称应力幅 / MPa & 426.9 & 597.7 & --- \\
计算寿命 / 次 & $3.0\times10^{4}$ & $6.9\times10^{3}$ & --- \\
200 次累积损伤 & $6.7\times10^{-3}$ & $2.9\times10^{-2}$ & $\le 1.0$ \\
疲劳安全系数 & 1.62 & 1.56 & $\ge 1.5$ \\
静强度安全系数 & 1.54 & 1.56 & $\ge 1.5$ \\
寿命裕度 & 约 150 & 约 34 & --- \\
\hline
\end{tabularx}
\end{table}

拉伸工况按 $R=0$ 代入计算流程：$\sigma_a=\sigma_m=302.1$\,MPa，$\sigma_{a,\mathrm{eq}}=302.1/(1-302.1/1\,033)=426.9$\,MPa，$N=(4\,512.7/426.9)^{1/0.2287}=3.0\times10^{4}$ 次（探针读数 30\,065 次），$D=200/3.0\times10^{4}=6.7\times10^{-3}$，$n_f=1/(302.1/929.7+302.1/1\,033)=1.62$，$n_s=930/604.1=1.54$。

抗扭工况无需作 Goodman 折算：$\sigma_m=0$，$\sigma_{a,\mathrm{eq}}=\sigma_a=597.7$\,MPa，$N=(4\,512.7/597.7)^{1/0.2287}=6.9\times10^{3}$ 次，$D=200/6.9\times10^{3}=2.9\times10^{-2}$。疲劳安全系数 $n_f=1/(597.65/929.7)=1.56$，与静强度安全系数 $n_s=930/597.65=1.56$ 完全相同。此现象并非计算巧合：完全反向时平均应力为零，$n_f=\sigma_e/\sigma_a$，而 4330V 在 $1\times10^{3}$ 次处的许用应力幅 929.7\,MPa 与屈服强度 930\,MPa 几乎相等，故 $n_f\approx\sigma_s/\sigma_a=n_s$。若材料改用屈服强度明显低于 $0.9\sigma_u$ 的档位，两者即会分开，$n_f$ 将低于 $n_s$。

两个工况均满足判据，但裕度并不宽裕：静强度安全系数 1.54 与 1.56 距判据 1.5 仅 3\%$\sim$4\%，抗扭累积损伤 $2.9\times10^{-2}$ 较同族其余工具高一个量级，寿命裕度约 34 倍。其原因为抗扭工况按完全反向取值，应力幅为拉伸工况的两倍，损伤约为其四倍。若按每趟 5 次行程计，$n=1\,000$ 次，抗扭累积损伤升至 $1.45\times10^{-1}$，仍小于判据 1.0，结论不变，但已进入需要关注的区间。

节点 774.29\,MPa 的读数是本节需要向设计方说明的一点：若以该值评定，静强度安全系数仅为 1.20，将不满足 1.5 的判据；而其与评定值相差 30\%，反映螺纹根部过渡圆角处的应力梯度很陡。本文按与轴向工况一致的规则取梯度平缓区，并建议下一版对该过渡圆角局部加密一档复核，以确认评定值的稳定性；同时应加大该处过渡圆角，避免因圆角偏小而使真实的应力集中升高。

控制部位为磁钢护套与中心管的连接螺纹根部。装配时应控制上扣扭矩，避免预紧过大在螺纹根部叠加拉应力；成品应按 NB/T\,47013.4 进行磁粉探伤，重点检查螺纹根部与台肩过渡处。
